\documentclass[12pt,a4paper]{ctexart}
\usepackage{geometry}
\usepackage{booktabs}
\usepackage[demo]{graphicx}
\usepackage{amsmath}
\geometry{left=2.5cm,right=2.5cm,top=2.5cm,bottom=2.5cm}
\title{AI 声明位置错误的反例}
\author{}
\begin{document}
\maketitle
\clearpage
\begin{abstract}
本样例用于验证：声明文书若被放到文献表之后是否会被判为位置错误。
摘要与关键词在此仅作占位。关键词：占位。
\end{abstract}
\newpage

\section{问题重述}
占位内容。
\section{问题分析}
占位内容。
\section{模型假设}
数据真实有效。
\section{符号说明}
符号约定如表 \ref{tab:s} 所示。
\begin{table}[htbp]
\centering
\caption{符号说明}
\label{tab:s}
\begin{tabular}{ll}
\toprule
符号 & 含义 \\
\midrule
$x$ & 决策变量 \\
\bottomrule
\end{tabular}
\end{table}
\section{模型的建立与求解}
目标函数见式 \eqref{eq:o}。
\begin{equation}
\min Z = \sum_{i} c_i x_i
\label{eq:o}
\end{equation}
\section{结果分析}
占位内容。
\section{模型检验}
\subsection{误差分析}
占位内容。
\subsection{灵敏度分析}
占位内容。
\subsection{稳健性分析}
占位内容。
\section{模型的优缺点与改进}
占位内容。
\section{结论}
占位内容。

\begin{thebibliography}{9}
\bibitem{ref1} 作者. 题名[J]. 刊名, 2020.
\end{thebibliography}

% ★★ 故意把 AI 声明放在参考文献之后 —— 这正是本反例要测的错误写法 ★★
\section*{AI 工具使用声明}
本参赛队在竞赛过程中使用了 AI 工具，主要用于语言润色，详细使用情况见支撑材料。
文中引用见 \cite{ref1}。

\appendix
\section{附录}
本论文没有用到程序。
\end{document}
